\section{怎样填写词语}

要想准确地填写词语，首先得在自己脑海里储存若干词语才行。俗话说：“巧媳妇难为无米之炊。”这是很容易明白的道理。可是在做词语填充题时，有的人明明学过某个词语却填不上来，或者错填成别的词语，这就说明，除了勤奋积累词汇之外，还得掌握填词的一些方法，提高分析、辨别、运用这些词语的能力，才能做好词语填充题。

那么，填写词语要注意些什么问题呢？

1. 注意辨析词类。

在填词时，我们务必分析要填的词在句子中的词性类别。如果词类不对，句子就不通顺，所填的词当然是错误的。例如：

（1）学好数理化与学好语文，二者不可偏（ ）。

（2）有些人对知识分子总有些不屑眼，他们这种偏（ ）不改变，是不利于社会主义建设事业的。

（3）我们要扫清前进道路上的（ ）碍。

（4）只能鼓励别人进步，不能（ ）碍别人进步。

第（1）句要填的是动词谓语的词素，第（2）句要填的是名词主语的词素，因此前者只能填“偏（废）”，后者只能填“偏（见）”，如果填颠倒了，整个句子就不通。第（3）句要填的是名词宾语的词素，第（4）句要填的是动词谓语的词素，因此前者只能填“障”碍”，后者只能填“阻”碍”，而不能倒置。
\newpage
根据这种辨析方法，我们就可以在确定所填词语具有什么词性的基础上选择恰当的字词进行填写了。例如：

（1）高空的白云和四周的雪峰清晰地倒（ ）水中，把湖光山色天影（ ）为晶莹的一体。在这（ ）静的湖上，唯一活动的东西就是天鹅。
（第一个要填的字，应该是“映”，与“倒”字组成动词谓语；第二个要填的字，应该是“融”，与“为”字组成动词谓语；第三个要填的字，应该是“幽”，与“静”字组成形容词定语。）

（2）秋天的阴山，象一座青铜的屏风安（ ）。在它们的北边，从阴山高处（ ）下来的深绿色的山坡，安（ ）地躺在黄河岸上，（ ）着阳光。这是多么平静的一个原野！

（第一个要填的字，应该是“放”，与“安”字组成动词谓语；第二个要填的字，应该是“拖”，与“下来”组成动补词组；第三个要填的字，应该是“闲”，与“安”字组成形容词状语；第四个要填的字，应该是“沐”，是动词谓语。）

2. 注意辨析词语之间的搭配。

因为要填的词语总是出现在句子中间，所以不注意前后词语之间的搭配关系是不行的。例如：

（1）文学作品主要是以人物形象来（ ）现社会生活的。

（2）今天的连云港又（ ）现出一派繁忙的景象。

（3）在“五讲四美三热爱”的热潮中（ ）现了大批先进人物。
\newpage
（4）我一想起他那亲切的面容，他那魁梧的身影就（ ）现在我的脑海里。

这四个句子中要填的字都是动词谓语的词素。第一句的动词是与宾语“生活”相搭配，只能用“（表）现”，第二句的动词是与宾语“景象”相搭配，只能用“（呈）现”，第三句的动词，前面有“在……热潮中”做状语，后面与宾语“人物”相搭配，因此只能用“（涌）现”；第四句的动词，后面有“在……脑海里”这个补语进行修饰和限制，因此只能填“（浮）现”。

根据这种辨析方法，我们就可以在弄清上下文词语的搭配关系的基础上选择恰当的字词进行填写了。例如：

（1）我们要以实际行动来（ ）护法律的权威，造成有法必（ ）、执法必（ ）、违法必（ ）的强大社会舆论。
（第一个要填的字，应该是“维”，与“护”字构成动词谓语，从而能与宾语“权威”恰当搭配；第二个要填的字，应该是“依”，从而与“法”字恰当搭配；第三个要填的字，应该是“严”，从而与“执法”恰当搭配；第四个要填的字，应该是“究”，从而与“违法”恰当搭配。）

（2）周恩来同志生前，从来不允许宣传他个人。恰恰在他（ ）后，全中国人民向自己的总理（ ）注了最珍贵的感情。“四人帮”这几个卑劣的小丑，（ ）想压制和玷污这种感情，结果他们自己被群众掀起的不可抗拒的怒潮（ ）没了。

（第一个要填的字，应该是“身”而不是“死”，因为从上下文来看，需要一个委婉的词相搭配；第二个要填的
\newpage
\noindent
字，应该是“倾”，与“注”字构成动词谓语，而能与宾语“感情”恰当搭配；第三个要填的字，应该是“安”，与“想”字构成动词谓语，以与前面的“小丑”和后面的“压制和玷污这种感情”所表达的贬义色彩相搭配；第四个要填的字，应该是“淹”，与“没”字构成动词谓语，以与前面的状语“被……怒潮”相搭配。）

3. 注意辨析上下文的词语结构方式。

如果上下文的句子是以并列方式出现的，那么这些句子中的词语往往具有相同的结构方式。这将给我们填词以有益的启示。例如：

（1）要想把事情想得（ ）密，话说得（ ）密，文章写得（ ）密，就必须自觉地掌握思维的形式和规律。

（2）柳亚子先生自辛亥革命以来，奔走呼号，声嘶力竭，举凡对于国势的兴（ ），民生的苦乐，友朋的（ ）散，乃至个人的（ ）处，无不一一发为诗歌。

我们从第一句的上下文看，要填的必定是几个同义词的词素，而想事情则要“（周）密”，说话得要“（严）密”，写文章当然要“（慎）密”了。再从第二句的上下文看，既然文句中出现了“民生的苦乐”（“苦乐”是由两个反义的词素构成的合成词），我们就可以判断：前后要填的字一定不是同义的词素，而是反义的词素。就“国势”而言，只能填“兴（衰）”；就“友朋”而言只能填“（聚）散”；就“个人”而言，只能填“（出）处”了。

根据这种辨析方法，我们就可以在弄清上下文的词语结构方式的基础上，选择恰当的字词进行填空了。例如：
\newpage
（1）江青为了篡党夺权，不管什么人，只要投其所好，（ ）其所为，粗腿（ ）得紧，马屁（ ）得响，“报告”（ ）得勤，谎话说得多，就委以重任。反之，谁反对她篡权，阻止她作恶，她就对你狠下毒手，（ ）加迫害，或把你（ ）起来，长期靠边，或把你放下去，终身劳改，或捏造罪名，置于死地。

（第一个要填的字，应该是“顺”，因为“顺其所为”与“投其所好”的词语结构方式相同，“顺”与“投”又是同义词。第二、三、四个要填的字，应该是“抱”、“拍”、“打”，因为这儿的三个词与“谎话说得多”的词语结构方式相同，而且“说”字是动词，所以相应要填的字也必然是动词，再考虑到这要填的三个字与前后词语的搭配关系，当然只能是“抱”粗腿、“拍”马屁和“打”报告了。第五个要填的字，应该是“横”；第六个要填的字，应该是“挂”，理由与上面所说大致相似。）

（2）寓言大多是借此喻（ ），借（ ）喻近，借古喻（ ），借小喻大。寓言的故事情节简（ ）浅显，却能说明复杂深（ ）的道理。

（第一、二、三个要填的字，应该是“彼”、“远”、“今”，因为这儿的三个词与“借小喻大”的词语结构方式相同，可以确定要填的字必然与前后相应的字意义相反。第四、五个要填的字，应该是“单”和“奥”，理由与上面所说大致相似。）

4. 注意辨析同义词的使用范围。

在填写词语的练习中，辨析同义词是重要的手段，而辨析同义词的使用范围，常常能帮助我们选择恰当的字词进行
\newpage
\noindent

填写。例如：

（1）今天，我校召开了“开创教育新（ ）面大会”，那种热烈的（ ）面使我兴奋极了。

（2）守卫边（ ）的解放军战士，夜以继日地在边（ ）线上巡逻。

第一句应填“（局）面”和“（场）面”，二者不能颠倒，因为“局面”是指一个时期内事情的状态，范围大，“场面”泛指一定场合下的情景，范围小。第二句要填的是“边（疆）”和“边（境）”，二者不能颠倒，因为“边疆”范围较大，“边境”范围较小。

5. 注意辨析同义词的轻重程度。例如：

（1）用词不妥贴，造句不合文法，行文缺乏条理，就会把意思弄得含混晦涩，令人（ ）解甚至（ ）解。

（2）我国文学源远流长，历史（ ）久，其中有许多列入世界文学之林而无愧的（ ）品。

第一句中要填的两个词的词义，肯定是前者轻，后者重，因为两词中间有表示递进关系的“甚至”一词。联系上文意思，第一个括号里应填“费”字，第二个括号里应填“误”字。“费解”只是难以理解，而“误解”则是错误地理解，程度就重多了。

第二句中要填的应是“（悠）久”和“（珍）品”。“悠久”是指年代久远，要比泛指时间长的“长久”程度重多了，何况“长久”与前面的“源远流长”的“长”字给人以用字重复之感。“珍品”是指珍贵的物品，而“佳品”、“妙品”等只是泛指好的物品罢了，词义的程度较“珍品”轻多了，何况只有“珍品”才与前面的定语“列入世界文学
\newpage
\noindent

之林而无愧”能够恰当地搭配。

注意辨析词义的感情色彩和风格特点。例如：
（1）当杨根思看好地形，仔细研究过作战计划后，就（ ）断地命令四班全体战士立即进行爆破。
（2）邓小平同志在人民大会堂（ ）见了金日成同志。

第一句只能填“果”字，而不能填“武”字，因为“果断”和“武断”虽然都有“作出决定，毫不犹豫”的意思，但是“果断”是褒义词，“武断”是贬义词，而从上文进行了调查研究的情况来看，这儿只能用褒义词。

第二句只能填“会”字，而不能填“看”字，更不能填“召”、“接”等字。因为“会见”一般用于新闻报道的文体，含有庄重的色彩，“看见”则用于一般文章，比较通俗，没有庄重的色彩；“召见”、“接见”虽然也用于新闻报道的文件和含有庄重的色彩，但是它们都用于上对下的场合，而邓小平同志和金日成同志是中、朝两国人民的领袖，他们之间的关系是平等的关系，因此只能用表示平等关系的“会见”。

以上所说的方法，虽然只用填写词素（或者写成“填写最恰当的字”）来说明，但是它适用的范围并不局限在填写词素上，也适用于填词、填成语、填词组。例如：
（1）我爱热闹，也爱（冷）（静）；爱群居，也爱（独）（处）。象今晚上，一个人在这苍茫的月光下，什么都可以想，什么都可以不想，便觉是个（自）（由）的人。
（朱自清《荷塘月色》）

（2）两年前的此时，即一九三一年的二月七日夜或八

\newpage
\noindent

日晨，是我们的五个青年作家同时遇害的时候。当时上海的报章都不敢载这件事，或者也许是（不）（愿）或（不）（屑）载这件事，只在《文艺新闻》上有一点隐约其辞的文章。
（鲁迅《为了忘却的记念》）

（3）风（卷）着雪花，狂暴地扫荡着山野、村庄，\\
（摇）（撼）着古树的躯干，（撞）（开）了人家的门窗，把破屋子上的茅草大把大把地（撕）下来向空中（扬）去，把冷森森的雪花（撒）进入家的屋子里，并且在光秃秃的树梢上怪声地怒吼着，（咆）（哮）着，仿佛世界上的一切，都是它的驯顺的奴隶，它可以任意地（践）（踏）他们，毁灭他们…… （峻青《党员登记表》）

我们只要在语文学习上不断下苦功，掌握丰富的词汇，并且能综合地灵活运用一些填写词语的方法，就一定能做好词语填充题。
